% ============================================================================
%   APPENDIX C — INDIVIDUAL RELIABILITY DIAGRAMS
% ============================================================================

\chapter{Individual Reliability Diagrams}
\label{app:reliability}

This appendix reproduces the individual reliability diagrams for each 
of the twelve configurations, in a compact grid format that facilitates 
visual comparison. Figure~\ref{fig:app-reliability-grid} presents the 
diagrams organized by model family (colors) and prompt variant (A/B in 
adjacent cells within each row).

Each panel plots the observed accuracy as a function of predicted 
confidence in ten equal-width bins. The dashed diagonal represents 
perfect calibration; points below the diagonal indicate overconfidence, 
points above indicate underconfidence. Configurations that exhibit 
mode collapse (Section~\ref{sec:res-mode-collapse}) produce degenerate 
reliability curves concentrated in a narrow range of the confidence 
axis, reflecting their inability to modulate confidence across cases.

\begin{figure}[htbp]
    \centering
    \includegraphics[width=\textwidth]{reliability_diagrams_grid.pdf}
    \caption{Individual reliability diagrams for the twelve 
        configurations. Colors indicate the model family: 
        \textcolor{red!80!black}{Mistral-Base} (red), 
        \textcolor{orange}{Mistral-Instruct} (orange), 
        \textcolor{blue}{Llama-Base} (blue), 
        \textcolor{black!70}{Llama-Instruct} (dark), 
        \textcolor{green!60!black}{Meditron} (green), 
        \textcolor{purple}{BioMistral} (purple). 
        The dashed diagonal represents perfect calibration. 
        Configurations with mode collapse show curves concentrated 
        in a narrow confidence range, unable to reflect the diversity 
        of clinical situations.}
    \label{fig:app-reliability-grid}
\end{figure}

Three visual patterns emerge from the grid. First, the two Mistral-Base 
panels (top row) show curves that lie systematically above the diagonal 
in the low-to-mid confidence range, indicating underconfidence: the 
model is more accurate than it claims. Second, the two Llama-Instruct 
panels show curves that lie systematically below the diagonal in the 
high-confidence range, indicating severe overconfidence: the model 
asserts wrong answers with high confidence. Third, the BioMistral and 
Llama-Base~B panels display near-vertical curves concentrated around a 
single confidence value, which is the visual signature of mode 
collapse: the model produces a nearly constant confidence estimate 
regardless of the input, and its accuracy on that narrow range is 
close to the base rate of the majority class.

The individual patterns visible in this grid complement the aggregate 
comparison of Figure~\ref{fig:res-reliability} of 
Chapter~\ref{chap:results} and support the finer interpretations 
provided in Sections~\ref{sec:res-h1} and~\ref{sec:res-mode-collapse}.